Tout d'abord, il faut trouvé une vitesse à l'équilibre. C'est à dire, la
vitesse à laquel le participant arrète de déaccéléré dans l'eau.
Ce point est important car c'est lui le centre de la linéarisation.
Hors, selon les coefficients données, la somme des forces est utilisé pour le
calcul. En effet, l'accélération peut être mis à zéro une fois l'équation isolé.
Par soucis du requis de 4 pages, les étapes d'isolations seront tous omises,
bien que tous disponibles en format propre.
\begin{align}
ma&=k_fmg+bv^2-mg \\
a&=g(k_f-1)+\frac{b}{m}v^2
\end{align}

Il est maintenant possible de trouvé la vitesse à l'équilibre:
\begin{align}
0&=g(k_f-1)+\frac{b}{m}v^2 \\
v_e&\approx0.9137\text{ m/s}
\end{align}

À partir de l'équation donnée par le guide étudiant, il est possible de
simplifié les vitesses et trouvé le terme à linéarisé:
\begin{align}
mv\frac{dv}{dt}&=mg(k_f-1)+bv^2 \\
\frac{dv}{dz}&=\frac{1}{v}g(k_f-1)+\frac{bv}{m}
\end{align}

La série de Taylor au premier degrée (avec v) est ceci:
\begin{align}
f(v)=f(v_e)+(\frac{\partial f}{\partial v})_e(v-v_e)
\end{align}

Hors, $f(v_e)$ à été calculé plus haut, et donnais $0$. Il peut donc être
enlevé. De plus, $(v-v_e)$ vas être remplacé par $\Delta v$. Il est
maintenant possible de faire la dérivé évalué à e;
\begin{align}
f(v_e)&=\frac{g(k_f-1)}{v_e}+\frac{bv}{m} \\
f(v_e)&=-\frac{g(k_f-1)}{v_e^2}+\frac{b}{m}
\end{align}

Mis dans l'équation de la série de Taylors;
\begin{align}
\frac{d\Delta v}{dz}=(-\frac{g(k_f-1)}{v_e^2}+\frac{b}{m})\Delta v
\end{align}

Afin de simplifié les prochaines démarches, une constante est mise pour évité
de dupliqué des grosses parties d'équations, de tel que:
\begin{align}
\lambda &= -\frac{g(k_f-1)}{v_e^2} \\
\lambda\Delta v&=\frac{d\Delta v}{dz}
\end{align}

Les termes semblable sont mis du même coté de l'équation:
\begin{align}
\frac{1}{\Delta v}d\Delta v=\lambda dz
\end{align}

les termes dérivées peuvent maintenant être intégrés afin d'obtenir un $z$:
\begin{align}
\int\frac{1}{\Delta v}d\Delta v&=\int\lambda dz \\
\ln|\Delta v|&=\lambda z + C
\end{align}

Le logarithme $\ln$ peut être enlevé en mettant les termes sur des $e$. 
\begin{align}
\Delta v=e^Ce^{\lambda z}
\end{align}

La constante d'intégration $C$ est juste une condition initial. En effet la
vitesses aux limites de l'intégral est connue. C'est la vitesse lors de
l'impact avec l'eau et de la vitesse à l'équilibre. Voici l'équation une fois
remplacé.
\begin{align}
\Delta v=\Delta v(0)e^{\lambda z}
\end{align}

Ne reste qu'a isolé $z$ de l'équation;
\begin{align}
\frac{1}{\lambda}\ln(\frac{\Delta v(z)}{\Delta v(0)})=z
\end{align}

Les deltas sont remplacé par leurs valeurs réelles, afin d'êtres mis dans MATLAB:
\begin{align}
\frac{1}{\lambda}\ln(\frac{v-v_e}{v_i-v_e})=z
\end{align}

La valeur de $z$ est recherché à $v=v_e*1.1\approx1$. et $v_i$ est la vitesse
lors de la rentrée dans l'eau. Cette dernière est obtenue de cette façon; ou la
hauteur est 10 mètres.
\begin{align}
E_p&=E_c \\
mgh_0&=\frac{1}{2}mv^2
\end{align}

Donc, $v_i\approx14$m/s. Une fois mis dans l'équation plus haut, on obtient une
profondeur de $z\approx4.2$ mètres.